%%%PAGE 275 rot=0 legibility=good summary=电磁感应中动量定理(安培力冲量与BLq、阻力正比于速度的冲量)、线框穿越磁场(q比、Q比)、连接体部分摩擦的摩擦力做功；页面下半为空白(仅有背面透字)
%%%BLOCK id=275-01 ch=12 sec=电磁感应·动量定理 kind=formula alt=07 cont=none y=0.03-0.12
\[
F_{\text{\unclear{安}}}\,t=BLq=\frac{B^2L^2x}{R} % CHECK: F后下标字被涂写(疑为“安”)；分子手写偏草，疑为 B^2L^2x
\]
\[
ft=kvt=kx
\]
%%%END
%%%BLOCK id=275-02 ch=12 sec=电磁感应·线框穿越磁场 kind=problem alt=07 cont=none y=0.14-0.29
% 图：线框(方框)速度 V_0 过竖直边界线，中间标 V_0/2，另一侧标“0 停止”
%FIG p275_a 0.15 0.135 0.35 0.215
\notefig[0.35]{figures/p275_a.png}

$q_1:q_2=1:1$

$Q_1:Q_2=3:1$ % 原稿 Q_1 处有涂写重叠，$Q_2$ 前有一多余竖笔
%%%END
%%%BLOCK id=275-03 ch=06 sec=功·摩擦力做功 kind=method alt=03 cont=none y=0.33-0.47
\textbf{连接体部分摩擦}

%FIG p275_b 0.125 0.37 0.30 0.43
\notefig[0.3]{figures/p275_b.png}
\[
W_f=\mu mg\,\frac L3
\]
%FIG p275_c 0.125 0.425 0.315 0.47
\notefig[0.3]{figures/p275_c.png}
\[
W_f=\mu mg\left(L+\frac L3\right)+\mu mg\,\frac L3
\]
%%%END
%%%PAGEEND 275
